%% ============================================================
%%  handout07.tex — Aula 7: Gradientes de política II e imitação
%%  Aprendizagem por Reforço — Universidade Federal de Uberlândia
%%  Prof. Marcelo Keese Albertini
%%  Adaptado e expandido de CS234 (Stanford), Emma Brunskill, Lecture 7
%%  Compilar com XeLaTeX (2x): latexmk -xelatex handout07.tex
%% ============================================================
\documentclass[11pt]{article}

\usepackage{handoutAprendizadoRL}
\usepackage{babel}
\babelprovide[import, main]{brazilian}
\usepackage{cancel}
\usepackage{array}

\setaula{Aula 7 --- PPO, GAE e imitação}
\setprof{Prof.\ Marcelo Keese Albertini}

%% ---- notação local (idêntica à dos slides da Aula 7)
\newcommand{\thv}{\theta}
\newcommand{\thk}{\thv_k}
\newcommand{\polth}{\pol_{\thv}}
\newcommand{\polk}{\pol_{k}}
\newcommand{\polnovo}{\pol'}
\newcommand{\gth}{\nabla_{\thv}}
\newcommand{\traj}{\tau}
\newcommand{\Vhat}{\hat{V}}
\newcommand{\Qhat}{\hat{Q}}
\newcommand{\Ahat}{\hat{A}}
\newcommand{\Adv}{A}
\newcommand{\Advpi}{A^{\pol}}
\newcommand{\Dkl}{D_{\mathrm{KL}}}
\newcommand{\dpi}{d^{\pol}}
\newcommand{\Lsub}[1]{\mathcal{L}_{#1}}
\DeclareMathOperator{\clip}{clip}
\DeclareMathOperator{\ent}{H}
\newcommand{\razao}{r_t(\thv)}
\newcommand{\Lclip}{\mathcal{L}^{\mathrm{CLIP}}}
\newcommand{\feat}{x}
\newcommand{\mufeat}{\mu}
\newcommand{\polE}{\pol^{*}}
\newcommand{\Dtot}{\mathcal{D}}
\newcommand{\wv}{w}

%% ---- respiro tipográfico para prosa em pt-BR com muitos termos técnicos
\setlength{\emergencystretch}{2.4em}
\hyphenation{es-pe-cia-lis-ta de-mons-tra-ções tra-je-tó-rias ex-po-nen-cial-men-te
  pro-ba-bi-li-da-des im-ple-men-ta-ção pa-ra-me-tri-za-da re-com-pen-sa
  Op-ti-mi-za-tion Re-in-force-ment}

%% ---- pseudocódigo numerado (mesmo estilo dos handouts anteriores)
\newcounter{plin}
\newcommand{\pl}[2][0em]{%
  \stepcounter{plin}%
  \par\noindent
  \hangindent=2.3em \hangafter=1
  \makebox[1.6em][r]{\color{rlCinza}\scriptsize\theplin:}\hspace{0.5em}\hspace*{#1}\ignorespaces#2\par}
\newenvironment{pseudo}%
  {\par\setcounter{plin}{0}\small\setlength{\parskip}{0.3ex}}%
  {\par}

\begin{document}

\capa{Aula 7 --- Gradientes de política II e Aprendizagem por imitação}
     {Amostragem por importância \textbullet\ Objetivo substituto e PPO \textbullet\ GAE \textbullet\
      Melhoria monotônica \textbullet\ Clonagem de comportamento, DAgger e RL inverso}
     {Prof.\ Marcelo Keese Albertini \textbullet\ Universidade Federal de Uberlândia}
     {Material adaptado e expandido de \textit{CS234: Reinforcement Learning}, Emma Brunskill
      (Stanford), Lecture 7. Leitura complementar: Schulman et al.\ (2016, 2017);
      Ross, Gordon e Bagnell (2011); Ziebart et al.\ (2008).}

\noindent
A Aula 5 nos deu um \emph{estimador} do gradiente da política. A Aula 6 mostrou que
saber a direção não basta: falta o \emph{tamanho do passo}, e medi-lo no espaço de
parâmetros é a coisa errada a fazer. Esta aula fecha o assunto de duas maneiras --- uma
prática (PPO e GAE, que é o que se implementa) e uma teórica (por que o esquema funciona)
--- e depois muda de assunto: e quando não há recompensa nenhuma, apenas alguém que já
sabe fazer a tarefa?

\begin{ideiabox}
  \textbf{Fio condutor da primeira metade.} Todo método desta aula responde à mesma
  pergunta: \emph{quanto posso mexer na política antes que os dados que tenho deixem de
  descrevê-la?} A resposta honesta é um limite teórico; a resposta que se usa na prática
  é o recorte do PPO. Vale conhecer as duas e saber onde uma vira a outra.

  \smallskip
  \textbf{Fio condutor da segunda metade.} Sem recompensa, sobra a demonstração. Ou se
  copia o \emph{comportamento} (clonagem, DAgger) ou se infere a \emph{intenção}
  (RL inverso). A primeira é barata e tem o especialista como teto; a segunda é cara e,
  em princípio, não tem.
\end{ideiabox}

\section*{Como usar este handout}
\addcontentsline{toc}{section}{Como usar este handout}

Este texto acompanha os slides da Aula 7 e é mais completo que eles em três pontos: as
passagens algébricas estão escritas por inteiro (em especial a soma geométrica do GAE e
a demonstração de melhoria monotônica), há pseudocódigo pronto para implementar, e há
exercícios ao final. As caixas \textcolor{rlAmbar}{âmbar} marcam perguntas para discussão
em aula --- não têm resposta escrita aqui de propósito.

\section{O problema do passo, e por que reponderar não resolve}

\subsection{De onde viemos}

O gradiente da política, na forma que já derivamos,
\[
  g \;=\; \gth J(\polth) \;=\;
  \E_{\traj\sim\polth}\Bigl[\sum_{t\ge0}\gamma^t\,\gth\log\polth(a_t\mid s_t)\,
  \Adv^{\polth}(s_t,a_t)\Bigr],
\]
é um estimador não-viesado, e sabemos reduzir sua variância (linha de base, vantagem,
estimadores de $n$ passos). Restam duas limitações, ambas sobre o \emph{uso} desse
gradiente:

\begin{enumerate}
  \item \textbf{Eficiência amostral ruim.} Cada lote de trajetórias é gasto em
    \emph{um} passo de gradiente e descartado, porque depois do passo a política mudou e
    os dados viraram off-policy.
  \item \textbf{Distância no espaço de parâmetros não é distância no espaço de
    políticas.} Este é o ponto conceitual. O espaço de políticas tabulares é o simplexo
    \[
      \Pi = \Bigl\{\, \pol \in \R^{\abs{\gS}\times\abs{\gA}} \;:\;
        \textstyle\sum_a \pol_{sa} = 1,\ \pol_{sa}\ge 0 \,\Bigr\},
    \]
    mas o gradiente anda em $\thv$. O mesmo $\norm{\Delta\thv}$ pode produzir uma
    mudança imperceptível de comportamento em uma região e uma política completamente
    diferente em outra. É por isso que ajustar o passo $\alpha$ é tão desagradável em
    gradiente de política e tão fácil em regressão logística.
\end{enumerate}

\begin{alertabox}
  A consequência prática do item 2 é assimétrica e vale enunciá-la: em aprendizado
  supervisionado, um passo grande demais produz uma época ruim e a próxima corrige. Em
  RL, um passo grande demais produz uma política ruim, que coleta \emph{dados ruins},
  com os quais se estima um gradiente ruim. O erro se realimenta.
\end{alertabox}

\subsection{Amostragem por importância: a tentativa óbvia}

Se o problema é que os dados vieram da política antiga, a ferramenta clássica para
corrigir é a amostragem por importância.

\begin{defbox}{Amostragem por importância}
  Para estimar $\E_{x\sim P}[f(x)]$ com amostras de $Q$ (com $Q(x)>0$ sempre que
  $P(x)>0$):
  \[
    \E_{x\sim P}[f(x)] \;=\; \E_{x\sim Q}\left[\frac{P(x)}{Q(x)}f(x)\right]
      \;\approx\; \frac{1}{\abs{\Dtot}}\sum_{x\in\Dtot}\frac{P(x)}{Q(x)}f(x),
      \qquad \Dtot\sim Q .
  \]
  A razão $P(x)/Q(x)$ é o \emph{peso de importância} de $x$.
\end{defbox}

O estimador é não-viesado para qualquer $Q$ admissível. O preço está na variância:
\begin{align*}
  \Var(\hat\mu_Q)
    &= \frac{1}{N}\Var_{x\sim Q}\!\left[\frac{P(x)}{Q(x)}f(x)\right] \\
    &= \frac{1}{N}\left(\E_{x\sim Q}\!\left[\frac{P(x)^2}{Q(x)^2}f(x)^2\right]
       - \E_{x\sim Q}\!\left[\frac{P(x)}{Q(x)}f(x)\right]^2\right)\\
    &= \frac{1}{N}\left(\mdestaque{rlLaranja}{\E_{x\sim P}\!\left[\frac{P(x)}{Q(x)}f(x)^2\right]}
       - \E_{x\sim P}[f(x)]^2\right).
\end{align*}
O termo destacado é o problema: se $P/Q$ for grande justamente onde $f^2$ é grande, a
variância explode. O viés continua zero --- e isso é uma armadilha, porque um estimador
não-viesado com variância enorme produz, em um único lote, números arbitrariamente
errados.

\subsection{Por que os pesos explodem em trajetórias}

Aplicando a ideia ao gradiente, com $\thv'$ a política que gerou os dados:
\begin{align*}
  g = \gth J(\thv)
   &= \E_{\traj\sim\thv}\Bigl[\sum_{t\ge0}\gamma^t \gth\log\polth(a_t\mid s_t)\Adv^{\thv}(s_t,a_t)\Bigr]\\
   &= \sum_{\traj}\sum_{t\ge0}\gamma^t P(\traj_t\mid\thv)\,\gth\log\polth(a_t\mid s_t)\Adv^{\thv}(s_t,a_t)\\
   &= \E_{\traj\sim\thv'}\Bigl[\sum_{t\ge0}
      \frac{P(\traj_t\mid\thv)}{P(\traj_t\mid\thv')}\,
      \gamma^t\gth\log\polth(a_t\mid s_t)\Adv^{\thv}(s_t,a_t)\Bigr].
\end{align*}

A razão de trajetórias colapsa exatamente como na Aula 5 --- a dinâmica desconhecida
e a distribuição inicial aparecem nos dois lados e se cancelam:
\[
  \frac{P(\traj_t\mid\thv)}{P(\traj_t\mid\thv')}
   = \frac{\cancel{\mu(s_0)}\prod_{t'=0}^{t}\cancel{P(s_{t'+1}\mid s_{t'},a_{t'})}\;\polth(a_{t'}\mid s_{t'})}
          {\cancel{\mu(s_0)}\prod_{t'=0}^{t}\cancel{P(s_{t'+1}\mid s_{t'},a_{t'})}\;\pol_{\thv'}(a_{t'}\mid s_{t'})}
   = \prod_{t'=0}^{t}\frac{\polth(a_{t'}\mid s_{t'})}{\pol_{\thv'}(a_{t'}\mid s_{t'})} .
\]

\begin{alertabox}
  O peso é um \textbf{produto de $t+1$ razões}. Ainda que cada fator esteja em
  $[0{,}95;\,1{,}05]$, com $t = 200$ o produto vive no intervalo
  $[0{,}95^{201};\,1{,}05^{201}] \approx [3\cdot10^{-5};\,2\cdot10^{4}]$. Diferenças
  minúsculas, multiplicadas duzentas vezes, viram uma diferença enorme. Na prática, um
  punhado de trajetórias domina o lote inteiro e as demais são numericamente ignoradas.
\end{alertabox}

A pergunta que organiza o resto da seção é a de Schulman e coautores: \emph{como
reaproveitar os dados da política antiga sem pagar o produto?}

\begin{discbox}
  O peso de importância explode com o horizonte. Isso sugere um remédio óbvio --- cortar
  o horizonte, ou descontar mais forte. Que preço se paga? Em que tipo de tarefa esse
  preço é aceitável e em que tipo é fatal?
\end{discbox}

\section{Objetivo substituto: uma razão em vez de um produto}

\subsection{A identidade exata de Kakade e Langford}

Existe uma expressão \emph{exata} para a diferença de desempenho entre duas políticas:

\begin{teobox}{Identidade de desempenho relativo (Kakade e Langford, 2002)}
  \[
    J(\polnovo) - J(\pol) \;=\; \frac{1}{1-\gamma}\,
      \E_{\substack{s\sim d^{\polnovo}\\ a\sim\polnovo}}\bigl[\Advpi(s,a)\bigr],
    \qquad
    \dpi(s) = (1-\gamma)\sum_{t\ge0}\gamma^t\,\Prob(s_t=s\mid\pol),
  \]
  onde $\dpi$ é a \emph{distribuição de ocupação descontada} de $\pol$.
\end{teobox}

A identidade tem uma leitura bonita: \emph{a nova política é melhor exatamente na medida
em que, nos estados que ela visita, ela escolhe ações com vantagem positiva segundo a
política antiga}. O incômodo é igualmente claro: os estados vêm de $d^{\polnovo}$, e
$\polnovo$ é justamente a política que ainda não executamos.

\subsection{A troca, e o que ela custa}

Trocamos $d^{\polnovo}$ por $\dpi$ --- que temos --- e reponderamos apenas a ação:

\begin{defbox}{Objetivo substituto $\Lsub{\pol}(\polnovo)$}
  \[
    \Lsub{\pol}(\polnovo) \defeq \frac{1}{1-\gamma}
      \E_{\substack{s\sim \dpi\\ a\sim\pol}}
      \left[\frac{\polnovo(a\mid s)}{\pol(a\mid s)}\;\Advpi(s,a)\right].
  \]
  Note a diferença essencial em relação à seção anterior: aqui há \textbf{uma razão por
  passo}, não um produto sobre a trajetória. É isso que torna o estimador utilizável.
\end{defbox}

A troca não é gratuita, e a teoria diz exatamente quanto custa:

\begin{teobox}{Limite de desempenho relativo (Achiam, Held, Tamar e Abbeel, 2017)}
  \[
    \Bigl| \bigl(J(\polnovo) - J(\pol)\bigr) - \Lsub{\pol}(\polnovo) \Bigr|
      \;\le\; C\,\sqrt{\E_{s\sim\dpi}\bigl[\Dkl(\polnovo\,\|\,\pol)[s]\bigr]},
    \qquad
    C = \frac{\gamma\,\max_{s}\abs*{\E_{a\sim\polnovo}[\Advpi(s,a)]}}{(1-\gamma)^2}.
  \]
\end{teobox}

Em palavras: o substituto é confiável na exata medida em que as duas políticas estão
próximas em divergência KL. A aproximação não é ruim --- ela é ruim \emph{longe}, com
erro controlado.

\begin{ideiabox}
  Repare no fator $(1-\gamma)^{-2}$ em $C$. Com $\gamma = 0{,}99$ ele vale $10^4$: a
  constante teórica é astronômica. Guarde isso; será o motivo pelo qual ninguém
  implementa o algoritmo teórico literalmente (Seção~4).
\end{ideiabox}

\section{PPO: duas maneiras de ficar perto}

O \emph{Proximal Policy Optimization} (PPO) é uma família de métodos que impõe a
restrição de KL \emph{de modo aproximado}, sem calcular gradientes naturais e sem
resolver o sistema linear do TRPO.

\subsection{Variante 1 --- penalidade KL adaptativa}

Resolve-se um problema irrestrito,
\[
  \thv_{k+1} = \argmax_{\thv}\;\Lsub{\thk}(\thv)
    - \beta_k\,\bar D_{\mathrm{KL}}(\thv\,\|\,\thk),
\]
e o coeficiente $\beta_k$ é reajustado entre iterações: se a KL medida no lote ficou
acima do alvo, aumenta-se $\beta$; se ficou muito abaixo, diminui-se. É um controlador
proporcional grosseiro sobre a restrição.

\subsection{Variante 2 --- objetivo recortado (a que se usa)}

\begin{defbox}{Objetivo recortado}
  Seja $\razao = \dfrac{\polth(a_t\mid s_t)}{\pol_{\thk}(a_t\mid s_t)}$ a razão de
  probabilidades \emph{de um passo}. Então
  \[
    \Lclip_{\thk}(\thv) = \E_{\traj\sim\polk}\Bigl[\sum_{t=0}^{T}
      \min\bigl(\razao\,\Ahat_t,\;\clip(\razao,\,1-\epsilon,\,1+\epsilon)\,\Ahat_t\bigr)\Bigr],
  \]
  com $\epsilon$ hiperparâmetro (tipicamente $0{,}2$), e
  $\thv_{k+1} = \argmax_{\thv}\Lclip_{\thk}(\thv)$.
\end{defbox}

A expressão parece opaca até que se separe em dois casos. Como $\Ahat_t$ é um número
fixo dentro do lote, o $\min$ escolhe entre duas retas em $r$:

\begin{center}
\begin{tikzpicture}[font=\small]
  % ---------- A > 0 ----------
  \begin{scope}
    \draw[->, rlCinza] (0,0) -- (5.0,0) node[below right, font=\scriptsize] {$r$};
    \draw[->, rlCinza] (0,-0.35) -- (0,2.15) node[left, font=\scriptsize] {$\Lclip$};
    \draw[rlAzul, very thick] (0,0) -- (3.3,1.55) -- (4.8,1.55);
    \draw[dashed, rlCinza!70] (2.2,-0.3) -- (2.2,1.03);
    \draw[dashed, rlCinza!70] (3.3,-0.3) -- (3.3,1.55);
    \node[below, font=\scriptsize, text=rlCinza] at (2.2,-0.3) {$1-\epsilon$};
    \node[below, font=\scriptsize, text=rlCinza] at (3.3,-0.3) {$1+\epsilon$};
    \fill[rlLaranja] (2.75,1.29) circle (2pt);
    \node[above left, font=\scriptsize, text=rlLaranja] at (2.75,1.29) {$r=1$};
    \node[font=\bfseries\small, text=rlAzul] at (2.4,2.45) {$\Adv > 0$ (ação boa)};
    \node[font=\scriptsize, text=rlCinza, align=center] at (2.4,-1.0)
      {sobe até $1+\epsilon$ e \emph{satura}:\\ sem incentivo a ir além};
  \end{scope}
  % ---------- A < 0 ----------
  \begin{scope}[xshift=7.6cm, yshift=1.4cm]
    \draw[->, rlCinza] (0,0) -- (5.0,0) node[above right, font=\scriptsize] {$r$};
    \draw[->, rlCinza] (0,-1.75) -- (0,0.45) node[left, font=\scriptsize] {$\Lclip$};
    \draw[rlAzul, very thick] (0,-1.03) -- (2.2,-1.03) -- (4.8,-2.25*0.78);
    \draw[dashed, rlCinza!70] (2.2,-1.75) -- (2.2,0);
    \draw[dashed, rlCinza!70] (3.3,-1.75) -- (3.3,0);
    \node[above, font=\scriptsize, text=rlCinza] at (2.2,0.02) {$1-\epsilon$};
    \node[above, font=\scriptsize, text=rlCinza] at (3.3,0.02) {$1+\epsilon$};
    \fill[rlLaranja] (2.75,-1.29) circle (2pt);
    \node[below right, font=\scriptsize, text=rlLaranja] at (2.75,-1.29) {$r=1$};
    \node[font=\bfseries\small, text=rlAzul] at (2.4,1.05) {$\Adv < 0$ (ação ruim)};
    \node[font=\scriptsize, text=rlCinza, align=center] at (2.4,-2.4)
      {constante abaixo de $1-\epsilon$;\\ \emph{decresce sem limite} acima};
  \end{scope}
\end{tikzpicture}
\end{center}

\medskip

\begin{itemize}
  \item \textbf{Se $\Adv_t > 0$} (a ação foi melhor que a média), queremos aumentar
    $r$. O objetivo cresce linearmente até $r = 1+\epsilon$ e então satura: o gradiente
    zera, e não há mais incentivo em empurrar a probabilidade da ação para cima.
  \item \textbf{Se $\Adv_t < 0$}, queremos diminuir $r$. O objetivo é \emph{constante}
    para $r < 1-\epsilon$ --- já reduzimos o suficiente, pare --- mas decresce sem
    limite do outro lado.
\end{itemize}

\begin{alertabox}
  Duas leituras que costumam escapar:
  \begin{itemize}
    \item O recorte remove o \emph{incentivo} de ir mais longe; ele \textbf{não impede}
      que $r$ fique grande. Se a primeira época de otimização já levou $r$ a $3$, o
      recorte não a traz de volta --- por isso PPO limita o número de épocas por lote e,
      em implementações sérias, monitora a KL e interrompe cedo.
    \item Para $\Adv<0$ o objetivo é ilimitado inferiormente: uma ação ruim pode ser
      punida arbitrariamente. É uma assimetria real do método, não um detalhe.
  \end{itemize}
\end{alertabox}

\begin{ideiabox}
  O $\min$ faz de $\Lclip$ um limitante \emph{pessimista} (inferior) do substituto não
  recortado: nunca somos otimistas a respeito do ganho de nos afastarmos da política que
  gerou os dados. É a mesma filosofia dos limites inferiores em RL conservador.
\end{ideiabox}

\subsection{O algoritmo, como se implementa}

\begin{algbox}{PPO com objetivo recortado}
\begin{pseudo}
\pl{Inicialize os parâmetros da política $\thv_0$ e do crítico $\wv_0$.}
\pl{\textbf{para} $k = 0, 1, 2, \dots$ \textbf{faça}}
\pl[1.2em]{Colete um lote $\Dtot_k$ de trajetórias executando $\pol_{\thk}$ por $T$ passos;
  guarde $\log\pol_{\thk}(a_t\mid s_t)$ para cada passo.}
\pl[1.2em]{Calcule os retornos e as vantagens $\Ahat_t$ por GAE truncado (Seção~\ref{sec:gae});
  normalize $\Ahat$ no lote.}
\pl[1.2em]{\textbf{para} época $= 1,\dots,K$ (tipicamente $K \in [3,10]$) \textbf{faça}}
\pl[2.4em]{\textbf{para} cada minilote de $\Dtot_k$ \textbf{faça}}
\pl[3.6em]{$\razao \leftarrow \exp\bigl(\log\polth(a_t\mid s_t) - \log\pol_{\thk}(a_t\mid s_t)\bigr)$}
\pl[3.6em]{Suba o gradiente de
  $\Lclip - c_1\,(\Vhat_{\wv}(s_t) - \hat R_t)^2 + c_2\,\ent[\polth(\cdot\mid s_t)]$.}
\pl[1.2em]{Interrompa as épocas cedo se $\hat D_{\mathrm{KL}}(\thv\|\thk)$ ultrapassar o alvo.}
\end{pseudo}
\end{algbox}

Três detalhes de implementação que separam um PPO que funciona de um que não funciona:

\begin{enumerate}
  \item \textbf{A razão é calculada por diferença de log-probabilidades}, nunca por
    divisão de probabilidades --- estabilidade numérica.
  \item \textbf{Normalizar $\Ahat$ por lote} (subtrair a média, dividir pelo desvio) é
    quase obrigatório; sem isso a escala do gradiente varia entre lotes e o
    $\epsilon$ fixo perde sentido.
  \item \textbf{O bônus de entropia} $c_2\,\ent[\polth]$ evita o colapso prematuro para
    uma política quase determinística, que mataria a exploração.
\end{enumerate}

\begin{alertabox}
  O ganho real de eficiência amostral está no laço interno (linha 5): várias atualizações
  de gradiente com o \emph{mesmo} lote. E é exatamente daí que vem a necessidade do
  recorte --- a partir da segunda época, os dados já são off-policy em relação à política
  que está sendo otimizada.
\end{alertabox}

\begin{discbox}
  O PPO tem quatro hiperparâmetros interagindo: $\epsilon$, o número de épocas $K$, o
  tamanho do lote e a taxa de aprendizado. Aumentar $K$ e diminuir $\epsilon$ tem efeito
  parecido sobre a KL final. São intercambiáveis? O que distingue um do outro?
\end{discbox}

\section{GAE: um botão único para o compromisso viés--variância}
\label{sec:gae}

Falta responder de onde vem o $\Ahat_t$ que aparece dentro do objetivo do PPO.

\subsection{Estimadores de $n$ passos e o erro TD}

Da Aula 5, os estimadores de vantagem com horizonte crescente:
\[
  \Ahat^{(1)}_t = r_t + \gamma V(s_{t+1}) - V(s_t), \qquad
  \Ahat^{(2)}_t = r_t + \gamma r_{t+1} + \gamma^2 V(s_{t+2}) - V(s_t), \quad\dots
\]
\[
  \Ahat^{(\infty)}_t = r_t + \gamma r_{t+1} + \gamma^2 r_{t+2} + \cdots - V(s_t)
    \;=\; G_t - V(s_t).
\]
Definindo o \textbf{erro TD}
\[
  \delta^V_t \defeq r_t + \gamma V(s_{t+1}) - V(s_t),
\]
todos eles se escrevem como somas telescópicas de erros TD:

\begin{align*}
  \Ahat^{(1)}_t &= \delta^V_t \\
  \Ahat^{(2)}_t &= \delta^V_t + \gamma\,\delta^V_{t+1} \\
  \Ahat^{(k)}_t &= \sum_{l=0}^{k-1}\gamma^l\,\delta^V_{t+l}
    \;=\; \sum_{l=0}^{k-1}\gamma^l r_{t+l} + \gamma^k V(s_{t+k}) - V(s_t).
\end{align*}

\begin{provabox}
Verificação do caso $k=2$, que já contém a ideia:
\begin{align*}
  \delta^V_t + \gamma\delta^V_{t+1}
  &= \bigl[r_t + \gamma V(s_{t+1}) - V(s_t)\bigr]
     + \gamma\bigl[r_{t+1} + \gamma V(s_{t+2}) - V(s_{t+1})\bigr]\\
  &= r_t + \gamma r_{t+1} + \gamma^2 V(s_{t+2}) - V(s_t)
     + \underbrace{\gamma V(s_{t+1}) - \gamma V(s_{t+1})}_{=0}.
\end{align*}
Os termos $\gamma^{l}V(s_{t+l})$ internos cancelam aos pares; sobram o
\emph{bootstrapping} na ponta ($+\gamma^k V(s_{t+k})$) e a linha de base no início
($-V(s_t)$). O caso geral segue por indução.
\end{provabox}

\subsection{A derivação do GAE}

\begin{defbox}{Estimador generalizado de vantagem (Schulman et al., 2016)}
  Média exponencialmente ponderada dos estimadores de $k$ passos, com peso
  $(1-\lambda)\lambda^{k-1}$ para $\Ahat^{(k)}$:
  \[
    \Ahat^{\mathrm{GAE}(\gamma,\lambda)}_t
      = (1-\lambda)\bigl(\Ahat^{(1)}_t + \lambda\Ahat^{(2)}_t + \lambda^2\Ahat^{(3)}_t + \cdots\bigr).
  \]
\end{defbox}

A forma acima é inútil para implementar --- exige todos os $\Ahat^{(k)}$. A mágica é que
ela colapsa. Escrevendo cada $\Ahat^{(k)}$ como soma de erros TD e reagrupando por
$\delta$:

\begin{align*}
  \Ahat^{\mathrm{GAE}}_t
  &= (1-\lambda)\Bigl(\delta^V_t
     + \lambda\bigl(\delta^V_t + \gamma\delta^V_{t+1}\bigr)
     + \lambda^2\bigl(\delta^V_t + \gamma\delta^V_{t+1} + \gamma^2\delta^V_{t+2}\bigr)
     + \cdots\Bigr)\\[0.2em]
  &= (1-\lambda)\Bigl(\delta^V_t\,(1+\lambda+\lambda^2+\cdots)
     + \gamma\delta^V_{t+1}\,(\lambda+\lambda^2+\cdots)
     + \gamma^2\delta^V_{t+2}\,(\lambda^2+\lambda^3+\cdots) + \cdots\Bigr)\\[0.2em]
  &= (1-\lambda)\left(\delta^V_t\,\frac{1}{1-\lambda}
     + \gamma\lambda\,\delta^V_{t+1}\,\frac{1}{1-\lambda}
     + \gamma^2\lambda^2\,\delta^V_{t+2}\,\frac{1}{1-\lambda} + \cdots\right)\\[0.2em]
  &= \mdestaque{rlLaranja}{\sum_{l=0}^{\infty}(\gamma\lambda)^{l}\,\delta^V_{t+l}} .
\end{align*}

\begin{ideiabox}
  O passo crítico é a segunda linha: trocar a ordem dos somatórios. O coeficiente que
  multiplica $\gamma^{l}\delta^V_{t+l}$ é $\lambda^{l}+\lambda^{l+1}+\cdots =
  \lambda^{l}/(1-\lambda)$, e o fator $(1-\lambda)$ da frente cancela exatamente esse
  denominador. Vale conferir o expoente com atenção: é $(\gamma\lambda)^l$, com
  \emph{ambos} os fatores.
\end{ideiabox}

O resultado é notável: uma média infinita de estimadores de todos os horizontes tem a
mesma forma --- e o mesmo custo --- de um único estimador de um passo, apenas com
$\gamma$ substituído por $\gamma\lambda$.

\subsection{Os dois extremos}

\begin{center}\small
\begin{tabular}{@{}>{\raggedright\arraybackslash}p{0.16\linewidth}
                 >{\raggedright\arraybackslash}p{0.36\linewidth}
                 >{\raggedright\arraybackslash}p{0.40\linewidth}@{}}
\toprule
 & $\lambda = 0$ & $\lambda = 1$ \\
\midrule
Estimador & $\Ahat_t = \delta^V_t$ & $\Ahat_t = \sum_{l\ge0}\gamma^l\delta^V_{t+l} = G_t - V(s_t)$\\
Equivale a & vantagem de TD(0) & vantagem de Monte Carlo \\
Variância & mínima (uma recompensa aleatória) & máxima (todas as recompensas) \\
Viés & alto sempre que $V\ne\Vpi$ & nenhum no gradiente resultante \\
\bottomrule
\end{tabular}
\end{center}

\begin{alertabox}
  Cuidado com a formulação de ``$\lambda=1$ é não-viesado''. Como \emph{estimador
  pontual} da vantagem, $\E[\Ahat_t] = \Qpi(s_t,a_t) - V(s_t)$, que só é igual a
  $\Advpi(s_t,a_t)$ se $V = \Vpi$. O que é não-viesado, para \emph{qualquer} $V$, é o
  \textbf{gradiente} resultante --- porque a diferença $\Vpi(s_t) - V(s_t)$ é uma função
  só do estado, e funções só do estado são linhas de base, cuja contribuição ao gradiente
  tem esperança zero (Aula 5).
\end{alertabox}

Na prática, escolhe-se $\lambda\in(0,1)$; valores entre $0{,}9$ e $0{,}97$ são o padrão
na literatura. O botão $\lambda$ é conceitualmente o mesmo dos traços de elegibilidade de
TD($\lambda$), aplicado à vantagem em vez do retorno.

\subsection{GAE truncado e a varredura reversa}

A soma infinita é impraticável; corta-se no fim do lote de $T$ passos:
\[
  \Ahat_t = \sum_{l=0}^{T-t-1}(\gamma\lambda)^l\,\delta^V_{t+l}.
\]

\begin{algbox}{Cálculo em uma varredura reversa, custo $\mathcal{O}(T)$}
\begin{pseudo}
\pl{$\Ahat_T \leftarrow 0$}
\pl{\textbf{para} $t = T-1$ \textbf{até} $0$ \textbf{faça}}
\pl[1.2em]{$\delta_t \leftarrow r_t + \gamma\,\Vhat(s_{t+1})\,(1-\mathrm{fim}_t) - \Vhat(s_t)$}
\pl[1.2em]{$\Ahat_t \leftarrow \delta_t + \gamma\lambda\,(1-\mathrm{fim}_t)\,\Ahat_{t+1}$}
\pl{Alvo do crítico: $\hat R_t \leftarrow \Ahat_t + \Vhat(s_t)$}
\end{pseudo}
\end{algbox}

O benefício é operacional: basta rodar a política por $T$ passos antes de atualizar, sem
esperar o episódio terminar. O truncamento introduz um viés adicional (a cauda além de
$T$ é ignorada), pequeno quando $(\gamma\lambda)^T \ll 1$. O indicador $\mathrm{fim}_t$
zera o \emph{bootstrapping} em estados terminais --- omiti-lo é o erro de implementação
mais comum nesta rotina, e produz um agente que ``espera recompensa depois da morte''.

\begin{discbox}
  Com $\gamma = 0{,}99$ e $\lambda = 0{,}95$, temos $\gamma\lambda \approx 0{,}94$ e a
  meia-vida do peso é de cerca de $11$ passos. Como escolher $T$ à luz disso? E o que
  acontece com o viés se o episódio típico durar $1000$ passos e $T = 128$?
\end{discbox}

\section{Por que isso funciona: melhoria monotônica}

\subsection{Do limite ao algoritmo}

Retirando o módulo do limite da Seção~2, fica a desigualdade que interessa:
\[
  J(\polnovo) - J(\pol) \;\ge\; \Lsub{\pol}(\polnovo)
     - C\sqrt{\E_{s\sim\dpi}\bigl[\Dkl(\polnovo\|\pol)[s]\bigr]}.
\]
O lado direito é um \textbf{minorante} do ganho verdadeiro --- e, ao contrário do lado
esquerdo, é estimável com amostras de $\pol$ apenas. Isso sugere um algoritmo: em vez de
maximizar o que não se pode medir, maximize o minorante.

\begin{teobox}{Proposição --- garantia de melhoria}
  Se $\polk$ e $\pol_{k+1}$ estão relacionadas por
  \[
    \pol_{k+1} = \argmax_{\polnovo}\;
      \Lsub{\polk}(\polnovo) - C\sqrt{\E_{s\sim d^{\polk}}\bigl[\Dkl(\polnovo\|\polk)[s]\bigr]},
  \]
  então $J(\pol_{k+1}) \ge J(\polk)$.
\end{teobox}

\begin{provabox}
A demonstração é curta e vale sabê-la de cor, porque o truque reaparece em vários lugares
do RL moderno. Basta exibir um ponto viável cujo valor do objetivo seja zero.

\begin{enumerate}
  \item $\polk$ pertence ao domínio da otimização, logo é um ponto viável.
  \item No ponto $\polnovo = \polk$, o substituto se anula:
    \[
      \Lsub{\polk}(\polk) = \frac{1}{1-\gamma}
        \E_{s\sim d^{\polk},\,a\sim\polk}\bigl[A^{\polk}(s,a)\bigr] = 0,
    \]
    porque $\E_{a\sim\pol}[\Advpi(s,a)] = 0$ para todo $s$ --- a vantagem tem média zero
    sob a própria política que a define.
  \item Também $\Dkl(\polk\|\polk)[s] = 0$ para todo $s$, de modo que o termo de
    penalidade se anula.
  \item Logo o objetivo avaliado em $\polk$ vale exatamente $0$; como $\pol_{k+1}$ é o
    \emph{máximo}, o valor ótimo é $\ge 0$.
  \item Pelo limite de desempenho relativo, $J(\pol_{k+1}) - J(\polk) \ge 0$.
    \hfill$\blacksquare$
\end{enumerate}

Duas observações sobre o alcance da demonstração. Primeira: ela não usa nada sobre a
forma de $\pol$, então continua valendo se restringirmos a otimização a uma classe
parametrizada $\Pi_\thv$ --- basta que $\polk\in\Pi_\thv$. É por isso que o argumento
sobrevive a redes neurais. Segunda: a melhoria garantida pode ser \emph{zero}; o teorema
diz que não pioramos, não que avançamos.
\end{provabox}

\begin{ideiabox}
  O esquema tem nome: \textbf{maximização de minorante} (\emph{majorize--minimize}).
  Constrói-se uma função que toca a verdadeira no ponto atual e fica sempre abaixo dela;
  maximiza-se essa função; repete-se. É o mesmo esqueleto do algoritmo EM e dos métodos
  de ponto proximal. Reconhecer o padrão é mais útil do que memorizar a fórmula.
\end{ideiabox}

\subsection{Por que ninguém usa isso literalmente}

A constante $C$ contém $(1-\gamma)^{-2}$. Com $\gamma = 0{,}99$, $C$ é da ordem de
$10^4\cdot\max_s\abs{\E_a[\Advpi]}$: a penalidade domina o substituto e o passo ótimo é
minúsculo. O algoritmo teoricamente correto é praticamente inútil.

\begin{center}\small
\begin{tabular}{@{}>{\raggedright\arraybackslash}p{0.24\linewidth}
                 >{\raggedright\arraybackslash}p{0.34\linewidth}
                 >{\raggedright\arraybackslash}p{0.34\linewidth}@{}}
\toprule
\textbf{Abordagem} & \textbf{O que faz com a KL} & \textbf{O que se perde} \\
\midrule
Teoria (acima) & penalidade com $C$ teórico & passos inúteis de tão pequenos \\
TRPO & \emph{restrição} $\bar D_{\mathrm{KL}} \le \delta$ resolvida por gradiente natural &
  custo por iteração (produto Hessiana--vetor, gradiente conjugado) \\
PPO & penalidade com $\beta$ ajustado, ou recorte da razão &
  a garantia formal; sobra uma heurística que funciona \\
\bottomrule
\end{tabular}
\end{center}

\begin{alertabox}
  Esta é uma lição metodológica que transcende o PPO: o valor do teorema não está na
  constante, está na \emph{forma} do limite. Ele diz que a penalidade certa é uma função
  crescente da KL --- e é essa forma que os métodos práticos preservam, trocando a
  constante por um hiperparâmetro calibrado empiricamente.
\end{alertabox}

\subsection{Balanço da primeira metade}

\begin{defbox}{O que fica de PPO e GAE}
\begin{itemize}
  \item \textbf{Eficiência amostral:} vários passos de gradiente por lote, em vez de um.
  \item \textbf{Estabilidade:} o recorte (ou a restrição de KL) aumenta a chance de
    melhoria monotônica, sem garanti-la.
  \item \textbf{Escopo:} converge para ótimos \emph{locais}. Como todo gradiente de
    política.
  \item \textbf{Relevância:} é o método padrão em controle contínuo e a base do
    ajuste por RLHF de modelos de linguagem.
\end{itemize}
\end{defbox}

\section{Aprendizagem por imitação: quando não há recompensa}

\subsection{O problema}

Em muitos domínios já \emph{existem} políticas de decisão muito boas --- executadas por
pessoas --- e o que se quer é automatizá-las. Há duas maneiras de aproveitar a pessoa:

\begin{center}\small
\begin{tabular}{@{}>{\raggedright\arraybackslash}p{0.30\linewidth}
                 >{\raggedright\arraybackslash}p{0.32\linewidth}
                 >{\raggedright\arraybackslash}p{0.32\linewidth}@{}}
\toprule
 & \textbf{Humano fornece recompensa} & \textbf{Humano demonstra} \\
\midrule
Canal & avalia decisões do agente & executa a tarefa \\
Vantagem & supervisão simples e barata por episódio & uma demonstração carrega muito mais informação que um número \\
Problema & complexidade amostral altíssima: milhões de episódios, cada um a ser avaliado &
  como generalizar para situações que o especialista nunca enfrentou? \\
\bottomrule
\end{tabular}
\end{center}

Há ainda um terceiro caminho, que é o das recompensas projetadas à mão
(\emph{reward shaping}). Recompensas densas no tempo guiam o agente de perto e aceleram
muito o aprendizado, mas escrevê-las é notoriamente frágil: o agente otimiza o que foi
escrito, não o que se queria. Uma penalidade por tempo mal calibrada produz um robô que
trava; um bônus por velocidade produz um robô que capota. Demonstrações são uma forma de
especificar a recompensa \emph{implicitamente}.

\begin{defbox}{Formulação do problema}
  \textbf{Dados:} espaços $\gS$ e $\gA$; modelo de transição $\tran{s'}{s,a}$;
  \textbf{nenhuma} função de recompensa; e um conjunto de demonstrações
  $(s_0,a_0,s_1,a_1,\dots)$ com ações extraídas da política do especialista $\polE$.

  \smallskip
  Três perguntas distintas, que dão origem a três famílias de métodos:
  \begin{enumerate}
    \item \textbf{Clonagem de comportamento:} dá para aprender $\polE$ diretamente por
      aprendizado supervisionado?
    \item \textbf{RL inverso:} dá para recuperar $\rec$?
    \item \textbf{Aprendizado por aprendizagem} (\emph{apprenticeship learning}): dá para
      usar o $\rec$ recuperado para obter uma boa política?
  \end{enumerate}
\end{defbox}

\begin{ideiabox}
  A escolha entre recompensa e demonstração é, no fundo, sobre \emph{qual canal de
  comunicação é mais fácil para o humano}. Especificar uma recompensa que produza direção
  defensiva é difícil; dirigir defensivamente por vinte minutos é fácil.
\end{ideiabox}

\subsection{Clonagem de comportamento}

A redução é direta: trate o problema como aprendizado supervisionado. Fixe uma classe de
políticas (rede neural, árvore de decisão) e ajuste $\polth$ aos pares
$(s_0,a_0), (s_1,a_1), \dots$ minimizando uma perda de classificação (ações discretas) ou
de regressão (ações contínuas).

Dois sucessos históricos merecem menção: ALVINN (Pomerleau, 1989), que dirigia um veículo
a partir de imagens de uma câmera com uma rede rasa, e o piloto de simulador de voo de
Sammut et al.\ (1992). Em versões modernas --- redes recorrentes ou \emph{transformers}
sobre janelas de observação, como no BC-RNN de Mandlekar et al.\ (2021) --- a clonagem é
extensivamente usada em manipulação robótica e funciona muito bem quando os dados cobrem
bem o espaço de estados visitado.

\begin{alertabox}
  Note o que a clonagem \textbf{não} exige: nenhum modelo de transição, nenhuma
  interação com o ambiente, nenhuma recompensa. É o método mais barato do curso inteiro.
  Todo o resto desta seção existe porque esse barateamento tem um preço.
\end{alertabox}

\subsection{O preço: erros compostos}

O aprendizado supervisionado supõe pares $(s,a)$ i.i.d.\ e ignora a estrutura temporal.
Se os erros fossem independentes no tempo, com probabilidade de erro $\le \varepsilon$
por passo, teríamos
\[
  \E[\text{número de erros}] \;\le\; \varepsilon T .
\]
Mas erros em um MDP não são independentes: um erro leva o agente a um estado que o
especialista nunca visitou, sobre o qual a política não tem informação, o que torna o
próximo erro mais provável, e assim por diante.

\begin{teobox}{Intuição do resultado de Ross e Bagnell}
  Se o agente erra pela primeira vez em $t$, todos os $T-t$ passos seguintes podem estar
  fora da distribuição de treino. Somando sobre o instante do primeiro erro:
  \[
    \E[\text{número de erros}] \;\lesssim\;
      \varepsilon\bigl(T + (T-1) + \cdots + 1\bigr) \;\propto\; \varepsilon T^2 .
  \]
  O enunciado formal é o Teorema 2.1 de Ross e Bagnell (AISTATS 2010); o argumento acima
  é uma aproximação da ideia, não a demonstração.
\end{teobox}

A diferença entre $\varepsilon T$ e $\varepsilon T^2$ é o que se chama de
\textbf{descasamento de distribuições}:

\begin{defbox}{A hipótese que quebra}
  No aprendizado supervisionado, $(x,y)\sim\Dtot$ tanto no treino quanto no teste. Em um
  MDP,
  \[
    \text{treino: } s_t \sim \Dtot_{\polE},
    \qquad
    \text{teste: } s_t \sim \Dtot_{\polth},
  \]
  e $\Dtot_{\polth}$ \emph{depende da política que estamos aprendendo}: a distribuição de
  teste é um resultado do treinamento, não um dado do problema.
\end{defbox}

\begin{ideiabox}
  É a patologia do RL \emph{off-policy} vista de outro ângulo: os dados vêm de uma
  política e a avaliação acontece sob outra. Lá a resposta foi reponderar (amostragem por
  importância); aqui a resposta será \emph{consertar os dados}.
\end{ideiabox}

\subsection{DAgger: agregação de conjuntos de dados}

\begin{algbox}{DAgger (Ross, Gordon e Bagnell, 2011)}
\begin{pseudo}
\pl{$\Dtot \leftarrow$ demonstrações iniciais do especialista; treine $\pol_1$ em $\Dtot$.}
\pl{\textbf{para} $i = 1,\dots,N$ \textbf{faça}}
\pl[1.2em]{Execute $\pol_i$ no ambiente e colete os estados visitados $\{s_1,\dots,s_m\}$.}
\pl[1.2em]{\textbf{Consulte o especialista} sobre cada um desses estados, obtendo
  $\{(s_j, \polE(s_j))\}$.}
\pl[1.2em]{$\Dtot \leftarrow \Dtot \cup \{(s_j, \polE(s_j))\}$ \quad (agregação, nunca descarte)}
\pl[1.2em]{Treine $\pol_{i+1}$ no $\Dtot$ acumulado.}
\end{pseudo}
\end{algbox}

A ideia em uma frase: \emph{obter rótulos do especialista ao longo do caminho que a
política aprendida realmente percorre}. O conjunto de treino passa a cobrir a
distribuição de estados induzida pela própria política, que é a distribuição de teste ---
e o descasamento desaparece. Sob as hipóteses do artigo, DAgger produz uma política
determinística estacionária com bom desempenho \emph{sob a distribuição que ela mesma
induz}, com erro voltando a escalar como $\varepsilon T$.

\begin{alertabox}
  \textbf{A limitação, que é séria.} O especialista precisa estar disponível
  \emph{durante} o treinamento, para rotular estados que ele não escolheu visitar --- e
  frequentemente estados ruins, em que ele nunca estaria. Duas dificuldades:
  \begin{itemize}
    \item O custo humano é recorrente, não um lote único de demonstrações.
    \item Rotular fora de contexto é difícil e ruidoso: ``o que você faria neste estado?''
      é uma pergunta muito mais difícil do que ``dirija bem por vinte minutos''. Em
      domínios com dinâmica rápida, a resposta pode nem ser bem definida.
  \end{itemize}
\end{alertabox}

\begin{discbox}
  Por que não basta simplesmente coletar mais demonstrações do especialista? Em que
  situação isso realmente resolveria o problema de erros compostos, e por que essa
  situação raramente ocorre na prática?
\end{discbox}

\section{RL inverso: inferir a intenção}

\subsection{Recompensa linear em atributos}

Suponha que a recompensa seja linear sobre um vetor de atributos do estado:
\[
  \rec(s) = \wv^{\!\top}\feat(s), \qquad \wv\in\R^n,\quad \feat:\gS\to\R^n .
\]
O valor de uma política herda a linearidade:
\begin{align*}
  V^{\pol}(s_0) &= \E_{\pol}\Bigl[\sum_{t\ge0}\gamma^t \rec(s_t)\,\Bigm|\, s_0\Bigr]
   = \E_{\pol}\Bigl[\sum_{t\ge0}\gamma^t\,\wv^{\!\top}\feat(s_t)\,\Bigm|\, s_0\Bigr]\\
   &= \wv^{\!\top}\,\underbrace{\E_{\pol}\Bigl[\sum_{t\ge0}\gamma^t \feat(s_t)\,\Bigm|\, s_0\Bigr]}_{\textstyle \mufeat(\pol)}
   \;=\; \wv^{\!\top}\mufeat(\pol).
\end{align*}

\begin{defbox}{Expectativas de atributos $\mufeat(\pol)$}
  $\mufeat(\pol)$ é a frequência descontada esperada com que cada atributo é encontrado
  ao seguir $\pol$ a partir de $s_0$. É o objeto que substitui a trajetória: duas
  políticas com o mesmo $\mufeat$ têm o mesmo valor sob \emph{qualquer} recompensa linear
  nesses atributos.
\end{defbox}

Se as demonstrações vêm de uma política ótima $\polE$ para a recompensa verdadeira
$\rec^* = \wv^{*\top}\feat$, então por definição de otimalidade
$V^{*} \ge V^{\pol}$ para toda $\pol$, ou seja
\[
  \wv^{*\top}\mufeat(\polE) \;\ge\; \wv^{*\top}\mufeat(\pol)
  \qquad \forall\,\pol .
\]
Identificar $\wv^*$ é, portanto, encontrar um vetor de pesos sob o qual o especialista
bata todo mundo.

\subsection{Ambiguidade: o problema estrutural do RL inverso}

\begin{alertabox}
  Não existe um $\rec$ único. Existe uma \textbf{infinidade} deles.
  \begin{itemize}
    \item $\wv = 0$ satisfaz a desigualdade acima com igualdade para toda política: se
      nada importa, o especialista é (vacuamente) ótimo.
    \item Escalar $\wv$ por qualquer constante positiva preserva a ordenação das
      políticas.
    \item Somar a $\rec$ um \emph{potencial} $\gamma\Phi(s') - \Phi(s)$ não altera a
      política ótima (\emph{reward shaping} de Ng, Harada e Russell, 1999).
    \item E, mais geralmente, há uma vizinhança inteira de recompensas para as quais a
      mesma política continua ótima.
  \end{itemize}
  Além disso, existem infinitas políticas \emph{estocásticas} que reproduzem as mesmas
  expectativas de atributos. O problema é duplamente subdeterminado, e todo método de RL
  inverso é, essencialmente, uma escolha de critério de desempate.
\end{alertabox}

\subsection{Casamento de expectativas de atributos}

Uma saída elegante: em vez de recuperar $\wv$, buscar diretamente uma política que
\emph{se comporte} como o especialista.

\begin{teobox}{Casamento de atributos (Abbeel e Ng, 2004)}
  Se $\norm{\mufeat(\pol) - \mufeat(\polE)}_1 \le \epsilon$, então para todo $\wv$ com
  $\norm{\wv}_\infty \le 1$ vale, pela desigualdade de Hölder,
  \[
    \abs{\wv^{\!\top}\mufeat(\pol) - \wv^{\!\top}\mufeat(\polE)} \le \epsilon .
  \]
\end{teobox}

\begin{provabox}
Hölder com o par conjugado $(1,\infty)$: para vetores $u,v$,
$\abs{u^{\!\top}v} \le \norm{u}_1\norm{v}_\infty$. Aplicando a
$u = \mufeat(\pol)-\mufeat(\polE)$ e $v = \wv$,
\[
  \abs{\wv^{\!\top}\bigl(\mufeat(\pol)-\mufeat(\polE)\bigr)}
    \le \norm{\mufeat(\pol)-\mufeat(\polE)}_1\,\norm{\wv}_\infty \le \epsilon\cdot 1 .
\]
Como $V^{\pol} = \wv^{\!\top}\mufeat(\pol)$, a desigualdade é exatamente uma cota sobre a
diferença de valores.
\end{provabox}

\begin{ideiabox}
  A garantia é forte de um jeito incomum: se as expectativas de atributos casam, a
  política é quase tão boa quanto a do especialista \textbf{para qualquer recompensa
  linear nesses atributos} --- inclusive a verdadeira, que continuamos sem conhecer.
  Contornamos a ambiguidade em vez de resolvê-la.
\end{ideiabox}

\begin{alertabox}
  O peso da hipótese migrou de $\rec$ para $\feat$. Se um atributo que importa ao
  especialista não estiver em $\feat$, nenhuma garantia o alcança --- a política pode
  casar todos os atributos observados e ainda assim se comportar de modo inaceitável na
  dimensão esquecida. Escolher $\feat$ é o trabalho difícil e raramente é discutido.
\end{alertabox}

\section{RL inverso de máxima entropia}

\subsection{Que distribuição escolher?}

Considere todas as trajetórias de $H$ passos em um MDP determinístico. Para uma recompensa
linear, uma política fica completamente especificada por sua distribuição sobre
trajetórias. Muitas distribuições reproduzem as expectativas de atributos observadas ---
qual escolher?

\begin{defbox}{Princípio da máxima entropia (Jaynes, 1957; Ziebart et al., 2008)}
  Escolha a distribuição \emph{sem nenhuma preferência adicional} além de casar as
  expectativas de atributos dos dados:
  \[
    \max_{P}\; -\sum_{\traj} P(\traj)\log P(\traj)
    \quad\text{sujeito a}\quad
    \sum_{\traj} P(\traj)\,\mufeat_{\traj} = \tilde\mufeat,
    \quad \sum_{\traj} P(\traj) = 1,
  \]
  onde $\mufeat_{\traj} = \sum_{s_i\in\traj}\feat(s_i)$ são as contagens de atributos de
  uma trajetória e $\tilde\mufeat = \frac{1}{m}\sum_{j=1}^{m}\mufeat_{\traj_j}$ é a média
  observada nas $m$ demonstrações.
\end{defbox}

\begin{alertabox}
  Atenção à mudança de definição: aqui $\mufeat_{\traj}$ é a soma de atributos ao longo de
  \emph{uma} trajetória, e não a expectativa descontada $\mufeat(\pol)$ da seção anterior.
  Os dois objetos usam a mesma letra na literatura e são frequentemente confundidos.
\end{alertabox}

\subsection{A forma exponencial e a verossimilhança}

A solução do problema de máxima entropia com restrições lineares é da família
exponencial:
\[
  P(\traj_j \mid \wv) = \frac{1}{Z(\wv)}\exp\bigl(\wv^{\!\top}\mufeat_{\traj_j}\bigr)
   = \frac{1}{Z(\wv)}\exp\Bigl(\sum_{s_i\in\traj_j}\wv^{\!\top}\feat(s_i)\Bigr),
  \qquad
  Z(\wv) = \sum_{\traj}\exp\bigl(\wv^{\!\top}\mufeat_{\traj}\bigr).
\]

A leitura é imediata, e é dela que vem o apelo do método: \textbf{caminhos de alto
retorno são exponencialmente mais prováveis, e caminhos de retorno igual são igualmente
prováveis}. O modelo não supõe alguém perfeito --- supõe alguém que erra, e que erra
menos quando o erro custa mais.

Em MDPs estocásticos a distribuição sobre caminhos depende também da dinâmica:
\[
  P(\traj_j \mid \wv, P) \;\approx\;
    \frac{\exp\bigl(\wv^{\!\top}\mufeat_{\traj_j}\bigr)}{Z(\wv,P)}
    \prod_{s_i,a_i\in\traj_j} P(s_{i+1}\mid s_i,a_i).
\]

\subsection{Ajustar \texorpdfstring{$\wv$}{w}: o gradiente é uma diferença de frequências}

Escolhe-se $\wv$ por máxima verossimilhança das demonstrações:
\[
  \wv^{*} = \argmax_{\wv} L(\wv) = \argmax_{\wv} \sum_{\text{exemplos}} \log P(\traj\mid\wv).
\]

\begin{teobox}{Gradiente da log-verossimilhança}
  \[
    \nabla L(\wv) \;=\; \tilde\mufeat \;-\; \sum_{\traj} P(\traj\mid\wv)\,\mufeat_{\traj}
      \;=\; \tilde\mufeat \;-\; \sum_{s_i} D(s_i)\,\feat(s_i),
  \]
  onde $D(s_i)$ é a \textbf{frequência esperada de visitação} do estado $s_i$ sob a
  política induzida pelos pesos atuais.
\end{teobox}

\begin{ideiabox}
  O gradiente é a diferença entre \emph{o que o especialista fez} e \emph{o que o modelo
  atual faria}. Se o modelo visita um estado mais do que o especialista, o peso dos
  atributos daquele estado desce. É a mesma estrutura do gradiente de modelos de
  máxima entropia em geral: dado observado menos expectativa do modelo.
\end{ideiabox}

O laço de otimização fica:

\begin{algbox}{MaxEnt IRL (Ziebart et al., 2008)}
\begin{pseudo}
\pl{Calcule $\tilde\mufeat$ a partir das demonstrações.}
\pl{Inicialize $\wv$.}
\pl{\textbf{repita}}
\pl[1.2em]{Resolva o MDP com $\rec_{\wv}(s) = \wv^{\!\top}\feat(s)$ (programação dinâmica
  \emph{soft}) para obter a política de máxima entropia.}
\pl[1.2em]{Propague para a frente as frequências de visitação $D(s)$ dessa política.}
\pl[1.2em]{$\wv \leftarrow \wv + \alpha\bigl(\tilde\mufeat - \sum_s D(s)\feat(s)\bigr)$.}
\pl{\textbf{até} convergir}
\end{pseudo}
\end{algbox}

\begin{alertabox}
  A linha 4 é o custo do método: \textbf{cada iteração resolve um MDP inteiro}. E a linha
  5 precisa do modelo de transição --- ou, na falta dele, da capacidade de interagir com
  o ambiente para amostrá-lo. Compare com a clonagem de comportamento, que não precisa de
  nenhum dos dois. Essa é a troca central desta segunda metade da aula.
\end{alertabox}

\subsection{De recompensas aprendidas a políticas}

O RL inverso entrega uma função de recompensa; o que normalmente se quer é uma
\emph{política} cujo desempenho iguale ou supere o do especialista. Duas rotas:

\begin{enumerate}
  \item \textbf{Em dois estágios:} aprenda $\rec$ por RL inverso e depois rode RL comum
    com essa recompensa. Simples, mas paga duas vezes o custo computacional e propaga os
    erros do primeiro estágio.
  \item \textbf{Diretamente:} aprenda a política sem materializar a recompensa. É o que
    faz o GAIL (Ho e Ermon, 2016), com uma formulação adversarial em que um discriminador
    tenta separar trajetórias do especialista das do agente, e o agente é treinado por
    gradiente de política para enganá-lo --- fechando o círculo com a primeira metade
    desta aula.
\end{enumerate}

\begin{ideiabox}
  Por que a recompensa é um objeto melhor que a política para transferir conhecimento?
  Porque ela é frequentemente mais \emph{compacta} e mais \emph{portável}: a intenção
  ``chegue rápido sem bater'' sobrevive à troca do carro, do mapa e da dinâmica; a
  política de direção, não. É também a razão pela qual o RL inverso pode, em princípio,
  superar o especialista --- se o especialista tinha boas intenções e execução ruim,
  otimizar a intenção supera imitar a execução.
\end{ideiabox}

\begin{discbox}
  A clonagem de comportamento, por construção, tem o especialista como teto. O RL
  inverso, não. O que exatamente permite ultrapassar quem nos ensinou --- e que hipótese
  precisa ser verdadeira para que isso aconteça de fato?
\end{discbox}

\section{Síntese e formulário}

\subsection{As cinco linhas que resumem a primeira metade}

\begin{center}\small
\begin{tabular}{@{}>{\raggedright\arraybackslash}p{0.30\linewidth}
                 >{\raggedright\arraybackslash}p{0.66\linewidth}@{}}
\toprule
Identidade exata &
  $J(\polnovo)-J(\pol) = \frac{1}{1-\gamma}\E_{s\sim d^{\polnovo}, a\sim\polnovo}[\Advpi(s,a)]$ \\[0.4em]
Substituto (estimável) &
  $\Lsub{\pol}(\polnovo) = \frac{1}{1-\gamma}\E_{s\sim\dpi,a\sim\pol}\bigl[\frac{\polnovo(a\mid s)}{\pol(a\mid s)}\Advpi(s,a)\bigr]$ \\[0.4em]
Preço da troca &
  $\bigl|(J(\polnovo)-J(\pol)) - \Lsub{\pol}(\polnovo)\bigr| \le C\sqrt{\E_{s\sim\dpi}[\Dkl(\polnovo\|\pol)]}$ \\[0.4em]
PPO recortado &
  $\Lclip = \E\bigl[\sum_t \min(\razao\Ahat_t,\ \clip(\razao,1-\epsilon,1+\epsilon)\Ahat_t)\bigr]$ \\[0.4em]
Vantagem (GAE) &
  $\Ahat_t = \sum_{l\ge0}(\gamma\lambda)^l\delta^V_{t+l}$, \quad
  $\delta^V_t = r_t + \gamma V(s_{t+1}) - V(s_t)$ \\
\bottomrule
\end{tabular}
\end{center}

\subsection{Mapa dos métodos de imitação}

\begin{center}\small
\begin{tabular}{@{}>{\raggedright\arraybackslash}p{0.185\linewidth}
                 >{\raggedright\arraybackslash}p{0.195\linewidth}
                 >{\raggedright\arraybackslash}p{0.115\linewidth}
                 >{\raggedright\arraybackslash}p{0.155\linewidth}
                 >{\raggedright\arraybackslash}p{0.225\linewidth}@{}}
\toprule
\textbf{Método} & \textbf{O que aprende} & \textbf{Precisa de $P$?} &
\textbf{Especialista no laço?} & \textbf{Limitação central} \\
\midrule
Clonagem & política, por supervisão & não & não & erros compostos ($\varepsilon T^2$) \\
DAgger & política, iterativamente & não & \textbf{sim} & custo humano recorrente \\
Casamento de atributos & política que casa $\mufeat$ & sim & não & depende inteiramente de $\feat$ \\
MaxEnt IRL & recompensa $\wv$ & sim & não & resolve um MDP por iteração \\
GAIL & política, de modo adversarial & não (usa simulador) & não & instabilidade do treino adversarial \\
\bottomrule
\end{tabular}
\end{center}

\subsection{O que você deve saber fazer}

\begin{defbox}{Capacidades esperadas}
\begin{itemize}
  \item Explicar por que os pesos de importância sobre trajetórias explodem, e por que a
    razão de \emph{um passo} do substituto não sofre do mesmo mal.
  \item Enunciar a identidade de Kakade--Langford e dizer qual termo é inestimável e por quê.
  \item Ler o objetivo recortado nos dois casos ($\Adv>0$ e $\Adv<0$) e explicar por que
    ele não impede que $r$ fique grande.
  \item \textbf{Derivar} a forma fechada do GAE a partir da média exponencialmente
    ponderada, indicando onde a soma geométrica é usada.
  \item Situar $\lambda = 0$ e $\lambda = 1$ no compromisso viés--variância, com a
    ressalva correta sobre o que é não-viesado.
  \item Reproduzir a demonstração de melhoria monotônica e apontar exatamente onde
    $\E_{a\sim\pol}[\Advpi] = 0$ é usada.
  \item Definir clonagem de comportamento e explicar em que ela difere de RL.
  \item Explicar o argumento $\varepsilon T$ contra $\varepsilon T^2$ e o que DAgger muda.
  \item Argumentar por que o RL inverso é subdeterminado, e como o casamento de atributos
    e a máxima entropia lidam com isso de maneiras diferentes.
\end{itemize}
\end{defbox}

\section{Exercícios}

\begin{exbox}{1 --- O produto que explode}
  Suponha $\polth(a_t\mid s_t)/\pol_{\thv'}(a_t\mid s_t) \in [1-\eta,\,1+\eta]$ para todo
  $t$, com $\eta = 0{,}02$.
  \begin{enumerate}
    \item Escreva o intervalo em que vive o peso de importância de uma trajetória de
      $t+1 = 100$ passos.
    \item Calcule numericamente os extremos. Compare com $\eta = 0{,}05$ e $t+1 = 500$.
    \item Explique por que o \emph{viés} do estimador continua zero enquanto isso.
  \end{enumerate}
  \textbf{Resposta parcial.} O intervalo é $[(1-\eta)^{100},(1+\eta)^{100}]
  \approx [0{,}133;\ 7{,}24]$: uma variação de $\pm2\%$ por passo vira um fator $54$ entre
  os extremos. Com $\eta=0{,}05$ e $500$ passos, $[(0{,}95)^{500},(1{,}05)^{500}]
  \approx [7{,}3\cdot10^{-12};\ 3{,}9\cdot10^{10}]$ --- quase vinte e duas ordens de grandeza.
\end{exbox}

\begin{exbox}{2 --- Os dois ramos do recorte}
  Escreva $\Lclip$ como função explícita de $r$ em cada caso, sem o $\min$:
  \begin{enumerate}
    \item Mostre que, para $\Adv > 0$, $\Lclip(r) = \Adv\cdot\min(r,\,1+\epsilon)$.
    \item Mostre que, para $\Adv < 0$, $\Lclip(r) = \Adv\cdot\max(r,\,1-\epsilon)$.
    \item Conclua que em ambos os casos $\Lclip(r) \le r\Adv$, isto é, o objetivo
      recortado é um limitante inferior do substituto.
  \end{enumerate}
\end{exbox}

\begin{exbox}{3 --- GAE em um episódio curto}
  Um episódio de três passos, com $\gamma = 0{,}9$:
  \[
    \begin{array}{c|ccc}
      t & 0 & 1 & 2\\\hline
      r_t & 0 & 0 & 10\\
      \Vhat(s_t) & 4 & 6 & 8
    \end{array}
    \qquad \Vhat(s_3) = 0 \ \ (\text{estado terminal}).
  \]
  \begin{enumerate}
    \item Calcule $\delta^V_0,\delta^V_1,\delta^V_2$.
    \item Calcule $\Ahat_0$ para $\lambda = 0$, $\lambda = 0{,}5$ e $\lambda = 1$.
    \item Verifique que, com $\lambda=1$, $\Ahat_0 = G_0 - \Vhat(s_0)$.
  \end{enumerate}
  \textbf{Resposta.} $\delta^V_2 = 10 + 0{,}9\cdot0 - 8 = 2$;
  $\delta^V_1 = 0 + 0{,}9\cdot8 - 6 = 1{,}2$;
  $\delta^V_0 = 0 + 0{,}9\cdot6 - 4 = 1{,}4$.
  Para $\lambda=0$: $\Ahat_0 = 1{,}4$.
  Para $\lambda=0{,}5$: $\Ahat_0 = 1{,}4 + 0{,}45\cdot1{,}2 + 0{,}45^2\cdot2 = 2{,}345$.
  Para $\lambda=1$: $\Ahat_0 = 1{,}4 + 0{,}9\cdot1{,}2 + 0{,}81\cdot2 = 4{,}1$, e de fato
  $G_0 = 0{,}9^2\cdot10 = 8{,}1$, com $8{,}1 - 4 = 4{,}1$.
\end{exbox}

\begin{exbox}{4 --- Onde a demonstração usa a hipótese}
  Na demonstração de melhoria monotônica, o passo 2 afirma
  $\E_{a\sim\pol}[\Advpi(s,a)] = 0$.
  \begin{enumerate}
    \item Demonstre essa identidade a partir de $\Advpi(s,a) = \Qpi(s,a) - \Vpi(s)$.
    \item Onde exatamente a demonstração quebraria se a vantagem fosse medida sob uma
      política diferente daquela que gera as ações?
    \item A garantia sobrevive se restringirmos $\polnovo$ a uma classe parametrizada
      $\Pi_\thv$? Justifique.
  \end{enumerate}
\end{exbox}

\begin{exbox}{5 --- $\varepsilon T$ contra $\varepsilon T^2$}
  Considere um agente clonado com erro por passo $\varepsilon = 0{,}01$ e horizonte
  $T = 200$.
  \begin{enumerate}
    \item Sob a hipótese i.i.d.\ ingênua, qual o número esperado de erros?
    \item Sob a estimativa $\varepsilon T^2$, qual é?
    \item Interprete: o que significa um número esperado de erros maior que $T$?
  \end{enumerate}
  \textbf{Comentário.} A resposta do item 3 é a chave: o limite $\varepsilon T^2$ passa a
  ser vacuamente satisfeito quando ultrapassa $T$, o que revela o regime em que a clonagem
  simplesmente não oferece garantia alguma.
\end{exbox}

\begin{exbox}{6 --- A ambiguidade, concretamente}
  Considere um MDP com dois estados $\{s_1,s_2\}$, duas ações e recompensa linear
  $\rec(s) = \wv^{\!\top}\feat(s)$ com $\feat(s_1) = (1,0)^{\!\top}$ e
  $\feat(s_2) = (0,1)^{\!\top}$.
  \begin{enumerate}
    \item Exiba dois vetores $\wv$ distintos e não proporcionais para os quais a mesma
      política determinística é ótima.
    \item Mostre que $\wv = 0$ torna \emph{qualquer} política ótima.
    \item Que restrição adicional o casamento de atributos impõe para eliminar a
      degenerescência do item 2?
  \end{enumerate}
\end{exbox}

\begin{exbox}{7 --- Implementação}
  Implemente PPO com GAE em um ambiente clássico de controle contínuo.
  \begin{enumerate}
    \item Registre, ao longo do treino, a KL empírica entre $\polth$ e $\pol_{\thk}$ ao
      fim de cada época interna. Ela cresce com $K$? Como?
    \item Varie $\epsilon \in \{0{,}05;\ 0{,}2;\ 0{,}5\}$ mantendo o resto fixo e
      compare as curvas de aprendizado e a fração de amostras recortadas.
    \item Varie $\lambda \in \{0;\ 0{,}95;\ 1\}$ e meça a variância empírica de $\Ahat$
      por lote. O comportamento bate com a tabela da Seção~\ref{sec:gae}?
    \item Remova o fator $(1-\mathrm{fim}_t)$ da varredura reversa e descreva o que
      acontece.
  \end{enumerate}
\end{exbox}

\begin{exbox}{8 --- Clonagem e DAgger no mesmo ambiente}
  Treine um especialista por RL até um bom desempenho e use-o como oráculo.
  \begin{enumerate}
    \item Faça clonagem de comportamento com $N$ demonstrações, para vários $N$, e trace
      o desempenho contra $N$.
    \item Implemente DAgger com o mesmo orçamento \emph{total} de rótulos e compare.
    \item Meça, em cada caso, a distância entre a distribuição de estados de treino e a
      de teste (por exemplo, com um classificador que tente distingui-las). A distância
      cai com DAgger?
  \end{enumerate}
\end{exbox}

\section*{Leituras}
\addcontentsline{toc}{section}{Leituras}

\begin{itemize}
  \item \textbf{Sutton e Barto (2018)}, cap.\ 13 --- gradientes de política; seção 12.1
    para a intuição de $\lambda$ nos traços de elegibilidade.
  \item \textbf{Schulman, Moritz, Levine, Jordan e Abbeel (2016)},
    \emph{High-Dimensional Continuous Control Using Generalized Advantage Estimation},
    ICLR --- o artigo do GAE; curto e vale cada página.
  \item \textbf{Schulman, Wolski, Dhariwal, Radford e Klimov (2017)},
    \emph{Proximal Policy Optimization Algorithms} --- o artigo do PPO.
  \item \textbf{Achiam, Held, Tamar e Abbeel (2017)},
    \emph{Constrained Policy Optimization}, ICML --- os limites de desempenho relativo
    usados aqui.
  \item \textbf{Ross, Gordon e Bagnell (2011)}, \emph{A Reduction of Imitation Learning
    and Structured Prediction to No-Regret Online Learning}, AISTATS --- DAgger e a
    análise de erros compostos.
  \item \textbf{Abbeel e Ng (2004)}, \emph{Apprenticeship Learning via Inverse
    Reinforcement Learning}, ICML --- casamento de expectativas de atributos.
  \item \textbf{Ziebart, Maas, Bagnell e Dey (2008)}, \emph{Maximum Entropy Inverse
    Reinforcement Learning}, AAAI.
  \item \textbf{Ho e Ermon (2016)}, \emph{Generative Adversarial Imitation Learning},
    NeurIPS.
  \item \textbf{Sun, Venkatraman, Gordon, Boots e Bagnell (2017)}, ICML --- a teoria na
    fronteira entre imitação e RL.
\end{itemize}

\vspace{0.6em}
\begin{ideiabox}
  \textbf{Próxima aula.} Aprender recompensas quando não há demonstrações completas, mas
  apenas \emph{preferências} entre pares de resultados ($y_1 \succ y_2$) --- o cenário dos
  \emph{dueling bandits} e o que sustenta o ajuste por feedback humano. E, na sequência, a
  exploração como problema de primeira classe.
\end{ideiabox}

\end{document}
